\documentclass[10pt,a4paper]{amsart}
\usepackage[margin=19mm]{geometry}
\usepackage{amsmath,amssymb,mathtools}
\usepackage[scr=rsfs,cal=euler]{mathalfa}
\usepackage{xfrac}
\usepackage[unicode,hidelinks,bookmarks=false]{hyperref}
\newcommand{\ie}{i.e.\ }
\newcommand{\cf}{cf.\ }
\newcommand{\resp}{respectively\ }
\newcommand{\Subsection}{\subsection}
\providecommand{\nobreakdash}{\nobreak\hbox{-}}
\setlength{\emergencystretch}{2em}
\multlinegap0pt
\usepackage{amssymb}
\usepackage{bm}
\usepackage[shortlabels]{enumitem}
\newtheorem{lemma}{Lemma}
\newcommand{\T}{{\mathbb T}}
\newcommand{\f}{\frac}
\newcommand{\mL}{\mathcal{L}}
\newcommand{\mR}{\mathcal{R}}
\title{Exact controllability of two-dimensional hydroelastic waves}
\author{Lizhe Wan, Jiaqi Yang}
\date{Source: arXiv:2603.01646v1 (CC BY 4.0)}
\begin{document}
\maketitle

\noindent\emph{Source and licence.} Exact controllability of two-dimensional hydroelastic waves. Version arXiv:2603.01646v1 (CC BY 4.0), distributed by arXiv under the Creative Commons Attribution 4.0 International licence, http://creativecommons.org/licenses/by/4.0/, as recorded in the arXiv licence metadata for this exact version and shown on the abstract page https://arxiv.org/abs/2603.01646v1. Full licence notice: \texttt{licenses/arxiv-2603.01646-CC-BY-4.0.txt}.\par\smallskip

\noindent\textit{Editorial scope.} Counted unit: the article's lemma t:CauchyEight (source0.tex lines 1162--1175). The article's complete original proof of it is delivered with it (source0.tex lines 1177--1223, one \texttt{proof} environment of 47 lines). No mathematics is written, repaired or re-derived: the statement and the proof are the article's own lines, in the article's own notation, and no retained line is substituted or re-flowed.  Retained uncounted as the source-local context the proof needs: lines 1153--1155.  Nothing was omitted from any retained span, so the package carries no omission record.  No retained line contains a citation, so the wrapper carries no bibliography and no bibliography entry.  No retained line contains a figure environment, an image-inclusion command or drawing code, so no image is delivered and the package declares no asset dependency.  Editorial formatting only, in addition: an AMS \texttt{amsart} preamble that restates the article's own declarations and macros used by the retained lines, namely the environments \texttt{lemma} on separate counters, together with the standard packages those lines need.  The retained environments print in this document as Lemma 1 in order of appearance, whereas the article prints the same items with the numbering of its own full document.  The complete native source of the article as distributed in the arXiv e-print for this exact version is bundled unchanged as \texttt{sources/195569/source0.tex}.
\begin{equation} \label{CauchyEight}
  (\mL_8 + \mR_8)\mathbf{h} = \mathbf{f}_8, \quad \mathbf{h}|_{\tau=0} = \mathbf{h}_{in}. 
\end{equation}

\begin{lemma} \label{t:CauchyEight}
Let $T>0$, and $s_1\geq 0$.
Suppose $\mathbf{h}_{in} \in H^{s_1}(\T)\times H^{s_1}(\T)$, $\mathbf{f}_8\in C([0,T];H^{s_1}(\T)\times H^{s_1}(\T))$.
There exists $\varepsilon_8>0$ depending on $T$ such that if 
\begin{equation*}
 \|\mR_8 \mathbf{f} \|_{C([0,T];H^{s_1}\times H^{s_1})} \lesssim \varepsilon_8 \|\mathbf{f}\|_{C([0,T];H^{s_1}\times H^{s_1})},
\end{equation*}
there exists a unique solution $\mathbf{h}\in C([0,T];H^{s_1}\times H^{s_1})$ of the Cauchy problem \eqref{CauchyEight} satisfying the estimate:
For $s\in [0, s_1]$, 
\begin{equation} \label{mR8EnergyEst}
 \| \mathbf{h}\|_{C([0,T]; H^s\times H^s)} \lesssim_s e^{CT}\left(\|\mathbf{h}_{in} \|_{H^s\times H^s} + \| \mathbf{f}_8\|_{C([0,T]; H^s\times H^s)}\right), 
\end{equation}
where the constant $C$ depends on $\varepsilon_8$.
\end{lemma}

\begin{proof}
We decompose $\mathbf{h} = \mathbf{h}_1+ \mathbf{h}_2$, where
\begin{equation} \label{mLEightCauchy}
\left\{
\begin{array}{lr}
\mL_8 \mathbf{h}_1 = \mathbf{f}_8 &  \\
\mathbf{h}_1|_{\tau=0} = \mathbf{h}_{in},&
\end{array}
\right. \quad
\left\{
\begin{array}{lr}
\mL_8 \mathbf{h}_2 = -\mR_8 \mathbf{h}_1 -\mR_8 \mathbf{h}_2 &  \\
\mathbf{h}_2|_{\tau=0} = \mathbf{0}.&
\end{array}
\right.
\end{equation}
The first system in \eqref{mLEightCauchy} is a linear system that can be solved explicitly.
There exists a unique solution $\mathbf{h}_1\in C([0,T]; H^s\times H^s)$, such that
\begin{equation}\label{es-h1}
 \| \mathbf{h}_1\|_{C([0,T]; H^s\times H^s)} \lesssim \|\mathbf{h}_{in} \|_{H^s\times H^s} + T\| \mathbf{f}_8\|_{C([0,T]; H^s\times H^s)}.
\end{equation}
We introduce the operator $\mathfrak{T}$ to denote the solution map of the first system in \eqref{mLEightCauchy}:
\begin{equation*}
\mathbf{h}_1:=\mathfrak{T}(\mathbf{f}_8,\mathbf{h}_{in}).
\end{equation*}
Thus, to solve the second system in \eqref{mLEightCauchy}, we need to prove that the following operator $\mathfrak{T}_1$ has a fixed point:
\begin{equation*}
\mathbf{w}= \mathfrak{T}_1(\mathbf{w}): =-\mathfrak{T}(\mathcal{R}_8\mathbf{h}_1+\mathcal{R}_8\mathbf{w},\mathbf{0}).
\end{equation*}
From \eqref{es-h1}, we have the uniform bound
\begin{align*}
&\|\mathfrak{T}_1(\mathbf{w})\|_{C([0,T]; H^s\times H^s)}\lesssim T\| \mathcal{R}_8(\mathbf{h}_1+\mathbf{w})\|_{C([0,T]; H^s\times H^s)}\\
\lesssim &T\|\mR_8\|_{C([0,T]; \mL (H^s\times H^s))}\left(\|\mathbf{h}_{in} \|_{H^s\times H^s} + T\| \mathbf{f}_8\|_{C([0,T]; H^s\times H^s)}+\| \mathbf{w}\|_{C([0,T]; H^s\times H^s)}\right), 
\end{align*}
and the contraction estimate
\begin{align*}
\|\mathfrak{T}_1(\mathbf{w}_1-\mathbf{w}_2)\|_{C([0,T]; H^s\times H^s)}&\lesssim T\| \mathcal{R}_8(\mathbf{w}_1-\mathbf{w}_2)\|_{C([0,T]; H^s\times H^s)}\\
&\lesssim T\|\mR_8\|_{C([0,T]; \mL (H^s\times H^s))}\| \mathbf{w}_1-\mathbf{w}_2\|_{C([0,T]; H^s\times H^s)}.
\end{align*}
From the above estimates,  if $\varepsilon_8$ is chosen such that $\varepsilon_8 T<\f12$, it follows that $\mathfrak{T}_1$ is a contraction. 
We then obtain the estimate from the contraction mapping principle:
\begin{align*}
 \| \mathbf{h}\|_{C([0,T]; H^s\times H^s)} \lesssim_s& (1+ T\|\mR_8\|_{C([0,T]; \mL (H^s\times H^s))})\|\mathbf{h}_{in} \|_{H^s\times H^s} + T\| \mathbf{f}_8\|_{C([0,T]; H^s\times H^s)} \\
+& T^2 \|\mR_8\|_{C([0,T]; \mL (H^s\times H^s))} \| \mathbf{f}_8\|_{C([0,T]; H^s\times H^s)}.
\end{align*}
This leads to the energy estimate \eqref{mR8EnergyEst}.
\end{proof}

\end{document}
